\chapter{Experimental Design and Benchmark Protocol}
\label{ch:experimental_design}

\section{Introduction}
\label{sec:exp_intro}
In empirical machine learning and computer vision, evaluation methodology dictates the validity of scientific conclusions. Flawed evaluation protocols—such as overlapping training and testing splits, threshold cherry-picking, or lenient true-positive matching criteria—can artificially inflate performance metrics, producing optimistic claims that collapse upon real-world deployment.

This chapter details the rigorous experimental design and benchmarking protocols established in TerraShift AI. We describe the academic OSCD benchmark protocol, formalize the mathematical framework of our **Strict Greedy 1-to-1 Event-Matching Algorithm**, explain our accounting for alert fragmentation and background false alarms, establish spatial false alert density metrics, and define the held-out generalization and latency measurement protocols.

\section{Academic Benchmark Protocol: OSCD}
\label{sec:exp_oscd}
To establish reproducibility against published academic literature, TerraShift AI is first evaluated against the **Onera Satellite Change Detection (OSCD)** benchmark \citep{daudt2018urban}. 

\subsection{Dataset Composition}
OSCD comprises 24 registered Sentinel-2 Level-1C/Level-2A image pairs acquired between 2015 and 2018 over major metropolitan regions worldwide (e.g., Paris, Montpellier, Beirut, Las Vegas, Dubai, Rio de Janeiro). The dataset is partitioned into:
\begin{itemize}
    \item \textbf{Training Set (14 image pairs):} Used for training supervised Siamese convolutional architectures.
    \item \textbf{Evaluation Test Set (10 image pairs):} Completely withheld during training; contains 3,077,936 evaluated ground-truth pixels.
\end{itemize}

\subsection{Pixel-Level Evaluation Metrics}
On OSCD, performance is measured at the raw pixel level by comparing predicted binary change masks $\hat{Y} \in \{0, 1\}^{H \times W}$ against ground-truth change annotations $Y \in \{0, 1\}^{H \times W}$. Let $TP_{\text{pix}}$, $FP_{\text{pix}}$, $TN_{\text{pix}}$, and $FN_{\text{pix}}$ denote pixel counts:
\begin{equation}
    \text{Precision}_{\text{pix}} = \frac{TP_{\text{pix}}}{TP_{\text{pix}} + FP_{\text{pix}}}
    \label{eq:pix_precision}
\end{equation}
\begin{equation}
    \text{Recall}_{\text{pix}} = \frac{TP_{\text{pix}}}{TP_{\text{pix}} + FN_{\text{pix}}}
    \label{eq:pix_recall}
\end{equation}
\begin{equation}
    \text{F1}_{\text{pix}} = \frac{2 \cdot \text{Precision}_{\text{pix}} \cdot \text{Recall}_{\text{pix}}}{\text{Precision}_{\text{pix}} + \text{Recall}_{\text{pix}}}
    \label{eq:pix_f1}
\end{equation}
\begin{equation}
    \text{IoU}_{\text{pix}} = \frac{TP_{\text{pix}}}{TP_{\text{pix}} + FP_{\text{pix}} + FN_{\text{pix}}}
    \label{eq:pix_iou}
\end{equation}
\begin{equation}
    \text{FPR}_{\text{pix}} = \frac{FP_{\text{pix}}}{FP_{\text{pix}} + TN_{\text{pix}}}
    \label{eq:pix_fpr}
\end{equation}

\section{The Polygon Fragmentation Problem}
\label{sec:exp_fragmentation_problem}
While pixel-level metrics are standard in academic segmentation challenges, they fail to reflect the operational realities of municipal surveillance and infrastructure monitoring. Planning authorities do not inspect individual $10\text{ m} \times 10\text{ m}$ pixels; they dispatch field inspectors to investigate discrete **vector polygons** representing physical land parcels.

Furthermore, naive object-level evaluation introduces a severe methodological flaw known as the **Polygon Fragmentation Artifact** (Figure \ref{fig:fragmentation_artifact}):
\begin{quote}
\textit{Suppose a single 50-hectare construction site is detected by a neural network not as a single coherent polygon, but as five fragmented sub-polygons due to internal boundary noise. Under naive overlap matching, each of the five sub-polygons intersects the ground-truth bounding box, and is therefore counted as a separate True Positive. Consequently, the detector is credited with $5\text{ TPs}$ for a single real event, artificially boosting precision while masking poor boundary cohesion.}
\end{quote}

\begin{figure}[htbp]
\centering
\fbox{
\begin{minipage}{0.9\textwidth}
\centering
\vspace*{0.2cm}
\textbf{The Polygon Fragmentation Dilemma}\\[0.2cm]
\begin{tabular}{|c|c|c|}
\hline
\textbf{Ground Truth Reference Event} & \textbf{Naive Overlap Matching} & \textbf{Strict Greedy 1-to-1 Matching} \\
\hline
Single 50-ha parcel & 5 fragmented alert polygons & 5 fragmented alert polygons \\
($N_{\text{ref}} = 1$) & all intersect reference box & sorted by confidence descending \\
\hline
\textbf{Result:} & $\mathbf{TP = 5}$, $FP = 0$ & $\mathbf{TP = 1}$ (Highest confidence) \\
& \textbf{Artificially Inflated Prec} & $\mathbf{FP_{\text{dup}} = 4}$ (Duplicate Detections) \\
& (Hides fragmentation) & \textbf{Strict Precision = 20.0\%} \\
\hline
\end{tabular}
\vspace*{0.2cm}
\end{minipage}
}
\caption{Comparison of naive overlap matching vs. strict greedy 1-to-1 matching.}
\label{fig:fragmentation_artifact}
\end{figure}

To guarantee absolute scientific honesty and prevent artificial metric inflation, TerraShift AI formalizes and enforces a **Strict Greedy 1-to-1 Event-Matching Protocol**.

\section{Localization-Aware Strict Greedy 1-to-1 Event-Matching Algorithm}
\label{sec:exp_greedy_matching}
Let $\mathcal{R} = \{r_1, r_2, \dots, r_K\}$ denote the set of coarse parcel-level reference change events in a corridor, each defined by a geographic bounding polygon $B(r_k)$ curated from multi-date Sentinel-2 visual and spectral inspection and known Saudi landmarks. Let $\mathcal{P} = \{p_1, p_2, \dots, p_M\}$ denote the set of predicted change polygons generated by a model after exact geometric clipping to the WGS84 corridor bounding box (\texttt{geo\_poly.intersection(shapely.box(*aoi\_bbox))}), with associated confidence scores $c(p_m) \in [0, 1]$.

To prevent sprawling unlocalized polygons from being credited as True Positives without imposing an arbitrary fixed percentage cap on genuine mega-project earthworks, a predicted polygon $p_m$ is a valid match candidate for reference parcel $r_k$ if and only if it satisfies both **physical parcel coverage** ($\text{RoC}$) and **localization purity** ($\text{PoC}$ and $\text{IoU}$):
\begin{equation}
    \text{RoC}(p_m, r_k) = \frac{\text{Area}(p_m \cap B(r_k))}{\text{Area}(B(r_k))}, \quad \text{PoC}(p_m, r_k) = \frac{\text{Area}(p_m \cap B(r_k))}{\text{Area}(p_m)}, \quad \text{IoU}(p_m, r_k) = \frac{\text{Area}(p_m \cap B(r_k))}{\text{Area}(p_m \cup B(r_k))}
    \label{eq:spatial_overlap}
\end{equation}
Under the standard benchmark protocol, $\text{Matches}(p_m, r_k)$ requires physical coverage ($\text{Area}(p_m \cap B(r_k)) \ge 3,500\text{ m}^2 \lor \text{RoC} \ge 0.08$) AND localization purity ($\left[\text{PoC} \ge 0.15 \land \text{IoU} \ge 0.05\right] \lor \text{PoC} \ge 0.30$). Any polygon that overlaps a reference parcel but fails the localization purity gate is penalized as a sprawling False Positive. We additionally evaluate all models across four localization regimes (\textit{Relaxed}, \textit{Standard}, \textit{Moderate}, and \textit{Strict}) to establish threshold sensitivity robustness.

Algorithm \ref{alg:greedy_matching} formalizes the strict greedy assignment procedure.

\begin{algorithm}[H]
\caption{Strict Greedy 1-to-1 Event-Matching Algorithm}
\label{alg:greedy_matching}
\begin{algorithmic}[1]
\REQUIRE Set of predicted polygons $\mathcal{P}$, Reference events $\mathcal{R}$, Negative zones $\mathcal{N}$
\STATE Sort $\mathcal{P}$ descending by confidence: $c(p_1) \ge c(p_2) \ge \dots \ge c(p_M)$
\STATE Initialize claimed reference set: $\mathcal{R}_{\text{claimed}} \leftarrow \emptyset$
\STATE Initialize counters: $TP \leftarrow 0$, $FP_{\text{dup}} \leftarrow 0$, $FP_{\text{bg}} \leftarrow 0$, $Hits_{\text{neg}} \leftarrow 0$
\FORALL{$p_m \in \mathcal{P}$}
    \STATE Find matching reference events: $\mathcal{M}_m = \{r_k \in \mathcal{R} \mid \text{Intersects}(p_m, r_k)\}$
    \IF{$\mathcal{M}_m \neq \emptyset$}
        \STATE Find unclaimed candidates: $\mathcal{U}_m = \mathcal{M}_m \setminus \mathcal{R}_{\text{claimed}}$
        \IF{$\mathcal{U}_m \neq \emptyset$}
            \STATE Select primary event: $r^* \leftarrow \text{first}(\mathcal{U}_m)$
            \STATE $\mathcal{R}_{\text{claimed}} \leftarrow \mathcal{R}_{\text{claimed}} \cup \{r^*\}$
            \STATE $TP \leftarrow TP + 1$ \COMMENT{Claimed as single True Positive}
        \ELSE
            \STATE $FP_{\text{dup}} \leftarrow FP_{\text{dup}} + 1$ \COMMENT{Overlaps already-claimed event}
        \ENDIF
    \ELSE
        \STATE $FP_{\text{bg}} \leftarrow FP_{\text{bg}} + 1$ \COMMENT{Unlinked background false alarm}
        \FORALL{$n_j \in \mathcal{N}$}
            \IF{$\text{Intersects}(p_m, n_j)$}
                \STATE $Hits_{\text{neg}} \leftarrow Hits_{\text{neg}} + 1$ \COMMENT{Negative control violation}
                \STATE \textbf{break}
            \ENDIF
        \ENDFOR
    \ENDIF
\ENDFOR
\RETURN $TP, FP_{\text{dup}}, FP_{\text{bg}}, Hits_{\text{neg}}, |\mathcal{R}_{\text{claimed}}|$
\end{algorithmic}
\end{algorithm}

\section{Event-Level Performance Metrics}
\label{sec:exp_metrics}
From the outputs of Algorithm \ref{alg:greedy_matching}, we compute three complementary precision metrics alongside recall and F1:

\begin{enumerate}
    \item \textbf{Strict 1-to-1 Event Precision ($\text{Prec}_{\text{strict}}$):} Penalizes both background false alarms and duplicate fragmented alerts. Represents the probability that any alert generated by the system corresponds to a unique, distinct real-world development project:
    \begin{equation}
        \text{Prec}_{\text{strict}} = \frac{TP}{TP + FP_{\text{dup}} + FP_{\text{bg}}} = \frac{TP}{|\mathcal{P}|}
        \label{eq:strict_precision}
    \end{equation}
    \item \textbf{Event Validity Precision ($\text{Prec}_{\text{validity}}$):} Measures the proportion of alerts that fall within genuine change zones, treating fragmented polygons within real projects as valid detections:
    \begin{equation}
        \text{Prec}_{\text{validity}} = \frac{TP}{TP + FP_{\text{bg}}}
        \label{eq:validity_precision}
    \end{equation}
    \item \textbf{Duplicate Alert Rate ($\text{Rate}_{\text{dup}}$):} Quantifies alert fragmentation:
    \begin{equation}
        \text{Rate}_{\text{dup}} = \frac{FP_{\text{dup}}}{|\mathcal{P}|}
        \label{eq:duplicate_rate}
    \end{equation}
    \item \textbf{Event Recall ($\text{Rec}_{\text{event}}$):} Measures the proportion of independent ground-truth projects successfully identified:
    \begin{equation}
        \text{Rec}_{\text{event}} = \frac{|\mathcal{R}_{\text{claimed}}|}{|\mathcal{R}|}
        \label{eq:event_recall}
    \end{equation}
    \item \textbf{Strict Event F1 Score ($\text{F1}_{\text{event}}$):} The harmonic mean of Strict Precision and Event Recall:
    \begin{equation}
        \text{F1}_{\text{event}} = \frac{2 \cdot \text{Prec}_{\text{strict}} \cdot \text{Rec}_{\text{event}}}{\text{Prec}_{\text{strict}} + \text{Rec}_{\text{event}}}
        \label{eq:event_f1}
    \end{equation}
\end{enumerate}

\section{Spatial False Alert Density Formulations}
\label{sec:exp_alert_density}
In operational GIS, raw false positive counts are meaningless without normalizing for the geographic area evaluated. A system generating 10 false alarms across $1,000\text{ km}^2$ is exceptionally reliable, whereas the same 10 false alarms in a $1\text{ km}^2$ neighborhood indicates failure.

We report two rigorous spatial density formulations:
\begin{enumerate}
    \item \textbf{Pooled False Alert Density ($D_{\text{pooled}}$):} Total false alarms across all corridors divided by total evaluated land area, normalized to a standard unit of $100\text{ km}^2$:
    \begin{equation}
        D_{\text{pooled}} = \frac{\sum_{k=1}^K \left(FP_{\text{dup}, k} + FP_{\text{bg}, k}\right)}{\sum_{k=1}^K \text{Area}_k} \times 100\text{ km}^2
        \label{eq:pooled_density}
    \end{equation}
    \item \textbf{Macro-Averaged False Alert Density ($D_{\text{macro}}$):} The unweighted arithmetic mean of per-corridor false alert densities:
    \begin{equation}
        D_{\text{macro}} = \frac{1}{K} \sum_{k=1}^K \left( \frac{FP_{\text{dup}, k} + FP_{\text{bg}, k}}{\text{Area}_k} \times 100\text{ km}^2 \right)
        \label{eq:macro_density}
    \end{equation}
\end{enumerate}

\section{Negative Zone Auditing and Generalization Testing}
\label{sec:exp_negative_zones}
\subsection{Negative Control Auditing}
Each evaluation corridor contains an explicitly designated **Negative Control Zone** (e.g., undisturbed limestone plateaus, natural sand dunes, stable sabkhas) verified by human remote sensing experts to have experienced zero anthropogenic or physical change across the observation epochs. Any detected polygon intersecting a negative control zone is recorded as an unambiguous algorithmic failure ($Hits_{\text{neg}}$). An operational production system must maintain $Hits_{\text{neg}} = 0$.

\subsection{Held-Out Corridor Generalization Protocol}
To evaluate geographical generalization, the **Al-Kharj Industrial & Logistics Corridor ($28.18\text{ km}^2$)** is designated as a strict held-out test corridor. All model parameters, detection thresholds ($\tau=0.48$ for Siamese, $\tau=0.12$ for AlphaEarth), and minimum area filters ($A_{\min}=3500\text{ m}^2$) were calibrated exclusively on the three core development corridors (Riyadh, Jeddah, Al-Jouf) and completely frozen prior to executing Al-Kharj evaluation. This provides an uncompromised assessment of transferability across unseen geographic regions.
\newpage
